% !TEX program = xelatex
% Example: Arabic journal article template
% Compile with: xelatex main.tex (twice for cross-references)
% Requires: polyglossia, bidi, Amiri font

\documentclass[12pt, a4paper]{article}

\usepackage{fontspec}
\usepackage{polyglossia}
\setmainlanguage[locale=mashriq]{arabic}
\setotherlanguage{english}
\newfontfamily\arabicfont[Script=Arabic]{Amiri}

\usepackage[margin=2.5cm]{geometry}
\usepackage{setspace}
\onehalfspacing
\usepackage{graphicx}
\usepackage{amsmath,amssymb}
\usepackage{booktabs}

\title{عنوانُ الورقةِ البحثية}
\author{اسمُ المؤلف \\
القسمُ، الجامعة \\
\LR{author@university.edu}}
\date{\today}

\begin{document}
\maketitle

\begin{abstract}
الملخّصُ يَشملُ المشكلةَ والمنهجَ والنتائجِ الرئيسةَ في 150 إلى 250 كلمة.
يَجبُ أنْ يُعطيَ القارئَ فكرةً كاملةً عنِ البحثِ دونَ قراءةِ الورقةِ كاملة.
\end{abstract}

\section{المقدمة}
تُعدُّ معالجةُ الإشاراتِ البيوميدية\LR{} مجالاً حيوياً في الهندسةِ الطبية.
تَهدفُ هذه الورقةُ إلى تقديمِ منهجيةٍ جديدةٍ لاستخراجِ الميزاتِ منْ إشاراتِ
تخطيطِ الدماغِ الكهربائيِّ \LR{(EEG)}.

\section{العملُ السابق}
تَناولتْ دراساتٌ سابقةٌ استخراجَ الميزاتِ بطرقٍ تقليديةٍ مثلَ تحليلِ
فورييه وتحليلِ المويجات. لكنَّ هذه الطرقَ تَعتمدُ على الخبرةِ اليدويةِ في
اختيارِ الميزاتِ المناسبة.

\section{المنهجية}
نَقترحُ منهجيةً تَجمعُ بينَ التحليلِ الزمني-التردديِّ والتعلّمِ العميق.
المعادلةُ الأساسيةُ للاستخراجِ هي:
\begin{equation}
\mathbf{f} = \text{CNN}(\text{CWT}(x(t)))
\end{equation}
حيثُ \(x(t)\) هيَ الإشارةُ الأصلية، و\(\text{CWT}\) هوَ تحويلُ المويجاتِ
المستمر، و\(\text{CNN}\) هيَ الشبكةُ العصبيةُ الالتفافية.

\section{النتائج}
يُوضّحُ الجدولُ \ref{tab:results} مقارنةً بينَ المنهجيةِ المُقترحةِ والطرقِ التقليدية.

\begin{table}[htbp]
\caption{مقارنةُ دقةِ التصنيفِ بينَ الطرقِ المختلفة.}
\label{tab:results}
\centering
\begin{tabular}{lc}
\toprule
\textbf{الطريقة} & \textbf{الدقة (\%)} \\
\midrule
تحليلُ فورييه & 78.5 \\
تحليلُ المويجات & 85.2 \\
المنهجيةُ المُقترحة & 93.7 \\
\bottomrule
\end{tabular}
\end{table}

\section{المناقشة}
تُظهرُ النتائجُ أنَّ المنهجيةَ المُقترحةَ تَتفوّقُ على الطرقِ التقليديةِ
بنسبةِ 8.5\%. هذا التحسّنُ يَعودُ إلى قدرةِ الشبكةِ العصبيةِ على تعلّمِ
الميزاتِ ذاتِ الصلةِ تلقائياً.

\section{الخاتمة}
قَدّمنا منهجيةً جديدةً لاستخراجِ الميزاتِ منْ إشاراتِ \LR{EEG} تَجمعُ
بينَ تحليلِ المويجاتِ والتعلّمِ العميق. النتائجُ تُظهرُ تحسّناً ملحوظاً
في دقةِ التصنيفِ مقارنةً بالطرقِ التقليدية.

\end{document}
